\documentclass[11pt,a4paper]{article}

\usepackage{fontspec}
\usepackage{amsmath,amssymb,amsthm,mathtools}
\usepackage[a4paper,margin=23mm]{geometry}

\usepackage{booktabs}
\usepackage{graphicx}
\usepackage{enumitem,needspace,placeins}
\usepackage{xcolor}
\definecolor{linkblue}{RGB}{24,56,88}
\usepackage[colorlinks=true,linkcolor=linkblue,citecolor=linkblue,urlcolor=linkblue,bookmarksnumbered=true]{hyperref}
\usepackage{bookmark}

\usepackage{polyglossia}
\setdefaultlanguage{hebrew}
\setotherlanguage{english}
\setmainfont{Liberation Serif}
\setmonofont{DejaVu Sans Mono}[Scale=.88]
\IfFileExists{fonts/NotoSansHebrew-Regular.ttf}{%
  \newfontfamily\hebrewfont[Path=fonts/,Extension=.ttf,UprightFont=*-Regular,BoldFont=*-Bold,ItalicFeatures={FakeSlant=0.18},ItalicFont=*-Regular,Script=Hebrew]{NotoSansHebrew}%
}{\newfontfamily\hebrewfont[Script=Hebrew,AutoFakeSlant=0.18]{Noto Sans Hebrew}}
\newcommand{\codepath}[1]{\textenglish{\path{#1}}}
\hypersetup{pdfcreator={XeLaTeX},pdflang={he}}
\AtBeginDocument{\renewcommand{\proofname}{הוכחה}\renewcommand{\abstractname}{תקציר}\renewcommand{\refname}{מקורות}\renewcommand{\contentsname}{תוכן העניינים}\renewcommand{\figurename}{איור}\renewcommand{\tablename}{טבלה}}

\newtheorem{theorem}{משפט}[section]
\newtheorem{lemma}[theorem]{למה}
\newtheorem{proposition}[theorem]{טענה}
\newtheorem{corollary}[theorem]{מסקנה}
\theoremstyle{definition}
\newtheorem{definition}[theorem]{הגדרה}
\theoremstyle{remark}
\newtheorem{remark}[theorem]{הערה}
\newtheorem{question}[theorem]{שאלה}
\newcommand{\Cspace}{\mathfrak C_{N,K}}
\newcommand{\R}{\mathbb R}
\newcommand{\Z}{\mathbb Z}
\newcommand{\TC}{\operatorname{TC}}
\newcommand{\tr}{\operatorname{tr}}
\newcommand{\id}{\operatorname{id}}
\newcommand{\doop}{\operatorname{do}}
\newcommand{\Lam}{\Lambda_{\mathrm{strict}}}
\newcommand{\Unif}{\operatorname{Unif}}
\newcommand{\Bern}{\operatorname{Bern}}
\newcommand{\relint}{\operatorname{relint}}
\newcommand{\argmax}{\operatorname*{arg\,max}}
\newcommand{\argmin}{\operatorname*{arg\,min}}
\setlist{itemsep=3pt,topsep=5pt}
\setlength{\emergencystretch}{2em}
\allowdisplaybreaks[2]
\raggedbottom

\newcommand{\EN}[1]{\textenglish{#1}}
\newcommand{\supp}{\operatorname{supp}}
\newcommand{\bits}{\{0,1\}}
\newcommand{\PP}{\mathsf{P}}
\newcommand{\NP}{\mathsf{NP}}
\newcommand{\UP}{\mathsf{UP}}
\newcommand{\BB}{\mathsf{BPP}}
\newcommand{\Atil}{\widetilde A}
\newcommand{\Prb}{\mathbb P}
\newcommand{\F}{\mathbb F}
\newcommand{\Eq}{\mathrel{:=}}
\newcommand{\gap}{2-\log_2 3}
\newcommand{\paperhead}[3]{\begin{center}\vspace*{2mm}{\LARGE\bfseries\color{linkblue}#1\par}\vspace{4mm}{\large #2\par}\vspace{2mm}{\small #3\par}\end{center}\vspace{3mm}}
\allowdisplaybreaks[1]

\newcommand{\RTLcell}[2]{\begin{minipage}[t]{#1}\raggedleft #2\end{minipage}}
\hypersetup{pdftitle={נקודות עיוורון מתורת המידע בביקורת סופית של קופסה שחורה},pdfkeywords={ביקורת סופית, שקילות תצפיתית, מתאם כולל, מינימקס, אורקלים, אלגבריזציה, Lean 4}}
\begin{document}
\paperhead{נקודות עיוורון מתורת המידע בביקורת סופית של קופסה שחורה}{אי־הבחנה מדויקת בפרופילים, חסמי מתאם והפרדת אורקלים שפורמלה בחלקה ב־\EN{Lean}}{קדם־פרסום מתוקן, גרסה \EN{2.0} \quad | \quad 25 בנובמבר 2026}
\begin{abstract}
אנו בוחנים מה יכול לאמת אוסף סופי וקבוע של בדיקות בינריות, כאשר האובייקט הנבדק מתואר באמצעות קואורדינטות התנהגותיות רבות שעשויות להיות תלויות זו בזו. לקונפיגורציה בת $N=2M$ ביטים המסומנת $B$, שההתפלגויות השוליות של הבלוקים בני שני הביטים שלה אחידות, לפרופיל ביקורת דטרמיניסטי $V=\alpha(B)\in\bits^m$ ולמתאם כולל $\TC(B)$, אנו מוכיחים שקיימת התפלגות על התומך המקורי המשמרת את מלוא החוק המשותף של $V$, ובה בעת מפחיתה את הסתברות ההצלחה של לפחות אחד מפרדיקטי ה־\EN{OR} על הבלוקים מ־$3/4$ לכל היותר ל־$2(m+\TC(B))/(N(2-\log_2 3))$. אנו משפרים חסם זה באמצעות אנטרופיה בינרית ומאפיינים אימות מדויק של משפחת פרדיקטים נתונה על ידי פירוק דרך פרופיל הביקורת. מודל נפרד של סדרות אינסופיות מציג עמימות סופית מתמשכת ופער מינימקס מרבי עבור מידות סיגמא־אדיטיביות. לבסוף, במודל האלגבריזציה האסימטרי של אהרונסון וויגדרסון, אנו בונים אורקל ישירות כך שההרחבה המולטילינארית היחידה שלו מאפשרת פעולת אגרגציה בשאילתה אחת, וכן
\[
\PP^A\subsetneq\UP^A=\NP^A\subseteq\PP^{\Atil[\mathrm{poly}]}.
\]
יחד עם אורקלי הקצה הידועים, מתקבלים כל שלושת הפרופילים האסימטריים המותרים בכל צילינדר תצפית סופי. הרכיבים האלגבריים ורכיבי האלכסון נתמכים במודולים של \EN{Lean 4.36.0} שעברו קומפילציה בנפרד; הוכחת השורש המאוחדת והתאמת מודל המכונה למודל הסטנדרטי נותרות מחוץ לקומפוזיציה שנבדקה. ההשלכות לביקורת מודלי שפה גדולים מותנות במודל התנהגותי מוגדר ובהתפלגות חלופית הניתנת למימוש. אין כאן טענה שבוצע ניסוי ב־\EN{LLM} או שנפתרה בעיית $\PP$ מול $\NP$.
\end{abstract}
\noindent\textbf{מילות מפתח:} ביקורת סופית; שקילות תצפיתית; מתאם כולל; מינימקס; סיבוכיות אורקלים; אלגבריזציה; \EN{Lean 4}.

\section{שאלה משותפת, שלושה מנגנונים נבדלים}
ביקורת סופית של קופסה שחורה מתעדת תשובות לאוסף נתון של שאילתות. הצלחה באוסף זה אינה מאמתת, ללא מבנה נוסף, כל תכונה שלא נצפתה. השאלה רלוונטית לאימות תוכנה ולביקורת מודלים; מחקרים אמפיריים על הצגה מכוונת של יכולת נמוכה יותר (\EN{sandbagging}) מספקים מוטיבציה, לא הוכחה למודל המופשט המשמש כאן \cite{taylor}. אנו שואלים איזה מידע הפרופיל משמר בפועל, אם התאימות הסופית ניתנת להרחבה גלובלית, ומה משתנה כשמשנים את פעולת השאילתה המותרת.

המאמר מפריד בין שלושה מנגנונים שונים מבחינה מתמטית. הראשון הוא אובדן מידע סופי: לשתי מידות יכול להיות אותו חוק מלא של הפרופיל הנצפה, אך הן עשויות להיבדל בפרדיקט היעד. השני הוא פער בכמתים: קיום עד לכל אוסף סופי של בדיקות אינו מבטיח קיום עד משותף. השלישי הוא גישה מוגבלת במשאבים: אותם נתוני יסוד עשויים לאפשר חישובים שונים כאשר הממשק מספק אגרגציה יקרה בעלות של שאילתה אחת. שלושת המנגנונים קשורים באמצעות ההבחנה בין אובייקט לבין הייצוג הנגיש שלו, אך איננו מניחים שאחת משלוש התוצאות נובעת מאחרת בלא הוכחה.

\begin{center}
\small\setlength{\tabcolsep}{7pt}
\begin{tabular}{lll}
\toprule
\RTLcell{0.29\linewidth}{\textbf{א. ביקורת סופית}} & \RTLcell{0.29\linewidth}{\textbf{ב. הרחבה גלובלית}} & \RTLcell{0.29\linewidth}{\textbf{ג. שינוי הגישה}}\\[3pt]
\midrule
\RTLcell{0.29\linewidth}{חוק פרופיל מדויק; תקציב אנטרופיה} & \RTLcell{0.29\linewidth}{צילינדרים סופיים; פער מינימקס} & \RTLcell{0.29\linewidth}{אותה טבלה בוליאנית; אגרגציה אלגברית}\\[3pt]
\RTLcell{0.29\linewidth}{$\alpha_*\mu=\alpha_*\nu$} & \RTLcell{0.29\linewidth}{$\forall Q\,\exists w_Q$ לעומת $\exists w\,\forall Q$} & \RTLcell{0.29\linewidth}{$\PP^A\subsetneq\UP^A=\NP^A$}\\[3pt]
\bottomrule
\end{tabular}
\end{center}
\noindent הדוגמאות אינן קובעות חסם תחתון אוניברסלי לכל שיטות הביקורת, ותוצאת האורקל אינה מחסום הוכחה חדש לבעיית המילניום ללא יחסוּת לאורקל \cite{cook}. יחסוּת ואלגבריזציה מגבילות מחלקות של שיטות הוכחה \cite{bgs,aw}; לעומתן, \EN{Natural Proofs} עוסקת בהבחנה מוגבלת־משאבים בטבלאות אמת הנתונות בשלמותן, תחת הנחה קריפטוגרפית \cite{rr}. אין לזהות בין שלושת המנגנונים רק משום שכולם משתמשים במונחים של אי־הבחנה.
\section{אי־הבחנה בפרופיל סופי ומתאם}
\subsection{אימות מדויק על תומך נתון}
יהי $S$ אוסף סופי של תצורות קבילות, תהי $\mu$ מידה שתומכה $S$, ותהי $\alpha:S\to\bits^m$ ביקורת דטרמיניסטית. נסמן $S_v=\{s\in S:\alpha(s)=v\}$ וכן $p_v=\mu(S_v)$. החלופות בסעיף זה הן מידות הסתברות שרירותיות על $S$; הדרישה שתומכיהן יהיו מוכלים ב־$S$ מחמירה יותר מהתרת תצורות חדשות שרירותיות. עבור יעד בוליאני $T:S\to\bits$, נגדיר
\[
q(T)=\sum_{v:S_v\ne\varnothing}p_v\min_{s\in S_v}T(s).
\]
\begin{proposition}[אי־הבחנה מדויקת ופירוק דרך הפרופיל]\label{prop:alias}
המינימום של $\mathbb E_\nu T$ על פני כל $\nu$ המקיימות $\alpha_*\nu=\alpha_*\mu$ שווה ל־$q(T)$, והוא מושג כאשר מעבירים את כל המסה של כל סיב לנציג הממזער את היעד. אף חלופה כזאת אינה יכולה להקטין את $\mathbb E_\mu T$ אם ורק אם $T$ קבוע על כל סיב שאינו ריק. עבור וקטור יעדים $T=(T_1,\ldots,T_M)$, כל היעדים ניתנים לאימות מדויק אם ורק אם $T=d\circ\alpha$ על $S$ עבור מפענח כלשהו $d$.
\end{proposition}
\begin{proof}
המסה של כל סיב קבועה. תרומתו היא לפחות מכפלת המסה שלו במינימום של $T$ עליו, וערך זה מושג באמצעות מסה נקודתית. מכיוון ש־$\mu$ מקצה מסה חיובית לכל איבר של $S$, סיב שעליו היעד אינו קבוע מאפשר ירידה ממש. החלת התוצאה הסקלרית על כל רכיב בנפרד נותנת את הפירוק.
\end{proof}
אם מותר להשתמש בפרדיקטים בינאריים דטרמיניסטיים שרירותיים כבדיקות, מספר הבדיקות המזערי לאימות מדויק של הווקטור הנתון על $S$ הוא $\lceil\log_2|T(S)|\rceil$: לדפוסי יעד שונים חייבים להיות פרופילים שונים, וקידוד בינארי חד־חד־ערכי של דפוסי היעד משיג את החסם. אם כל $2^M$ הדפוסים של \EN{OR} מופיעים, המספר הוא בדיוק $M$. הטענה מתייחסת למשפחת יעדים נתונה, ולא לכל התכונות האפשריות של מערכת מבוקרת.

\subsection{חסם כמותי לבלוקים מתואמים}
יהי $B=(B_1,\ldots,B_M)$ כאשר $B_i\in\bits^2$, מתקיים $N=2M$, וההתפלגות השולית של כל $B_i$ אחידה על פני ארבעת הזוגות. תהי $\mu$ ההתפלגות המשותפת שלו, ונסמן $S=\supp\mu$ וכן $V=\alpha(B)\in\bits^m$. כל האנטרופיות מחושבות בבסיס שתיים. נגדיר את המתאם הכולל \cite{watanabe,cover} ואת יעד הבלוק כך:
\[
\TC(B)=\sum_{i=1}^M H(B_i)-H(B),\qquad
T_i(B)=B_{i,1}\lor B_{i,2}.
\]
מכאן $\Prb_\mu[T_i=1]=3/4$. אין מניחים אי־תלות של הבלוקים או של הבדיקות. נגדיר את המתאם הכולל המותנה באמצעות $\TC(B\mid V)=\sum_iH(B_i\mid V)-H(B\mid V)\ge0$. מכיוון ש־$V$ דטרמיניסטי,
\begin{equation}\label{eq:budget}
\sum_{i=1}^M I(B_i;V)=H(V)+\TC(B)-\TC(B\mid V)\le m+\TC(B).
\end{equation}
\begin{theorem}[אי־הבחנה מדויקת בפרופיל תוך התחשבות במתאם]\label{thm:alias}
קיימים אינדקס $i$ ומידת הסתברות $\nu$ על $S$, כך ש־$\alpha_*\nu=\alpha_*\mu$ ומתקיים
\begin{equation}\label{eq:linear}
\Prb_\nu[T_i=1]\le\min\left\{\frac34,\frac{2(m+\TC(B))}{N(\gap)}\right\}.
\end{equation}
אותה טענה מתקיימת כאשר מחליפים את $m$ ב־$H(V)$.
\end{theorem}
\begin{proof}
לפי \eqref{eq:budget}, קיים $i$ שעבורו $I(B_i;V)\le (m+\TC(B))/M$. בכל סיב לא ריק, בוחרים נציג שמקיים $B_i=00$ אם קיים נציג כזה ב־$S$, ונציג שרירותי אחרת. מקצים לו מסה $p_v$. כך מתקבלת $\nu$ שהתפלגות הפרופיל שלה זהה בדיוק להתפלגות המקורית. תהי $E_i$ קבוצת הפרופילים שבסיביהם בתומך לא מופיע $00$ בבלוק $i$, ונסמן $q_i=\Prb[V\in E_i]$. אז $\Prb_\nu[T_i=1]=q_i$, וכן $q_i\le3/4$. מכיוון שאנטרופיית הבלוק המותנית היא לכל היותר $\log_2 3$ על $E_i$ ולכל היותר $2$ מחוצה לה, מתקיים $I(B_i;V)=2-H(B_i\mid V)\ge q_i(2-\log_2 3)$. הצבה נותנת את החסם.
\end{proof}
כדי להבטיח כמותית ש־$\mathbb E_\nu T_i\ge\varepsilon$ לכל $i$ ולכל חלופה על $S$ המשמרת את הפרופיל, תנאי הכרחי הוא $m+\TC(B)\ge M\varepsilon(2-\log_2 3)$. בפרט, החסם הליניארי המקורי נותן פער כלפי מטה שאינו נעלם כאשר $m=o(N)$ וכן $\limsup\TC(B)/N < \tfrac38(2-\log_2 3)\approx0.1556$. כדי שהחסם ידעך ל־$o(1)$ נדרש $m+\TC(B)=o(N)$; עצם קיומו של יחס מתאם חיובי וקבוע מתחת לסף אינו גורר דעיכה כזאת.

\subsection{סף אנטרופי משופר}
הבנייה שאינה חורגת מן התומך מאפשרת חסם חד יותר בלי לשנות את הביקורת. נסמן ב־$h_2$ את האנטרופיה הבינארית ונגדיר, עבור $0\le q\le3/4$,
\begin{equation}\label{eq:phi}
\phi(q)=h_2(1/4)-(1-q)h_2\!\left(\frac{1}{4(1-q)}\right).
\end{equation}
פונקציה זו רציפה ועולה ממש, ומקיימת $\phi(0)=0$ וכן $\phi(3/4)=h_2(1/4)$. בפנים התחום נגזרתה היא $-\log_2(1-1/[4(1-q)])>0$.
\begin{proposition}[חיזוק לא ליניארי]\label{prop:nonlinear}
אותו אינדקס ואותה בנייה ממשפט \ref{thm:alias} מקיימים
\[
\phi(q_i)\le\frac{H(V)+\TC(B)}{M}.
\]
לפיכך, לכל $\varepsilon\in(0,3/4]$ קבוע, כדי לאמת כל יעד ברמה $\mathbb E_\nu T_i\ge\varepsilon$ כנגד כל החלופות המשמרות את הפרופיל, נדרש $m+\TC(B)\ge M\phi(\varepsilon)$.
\end{proposition}
\begin{proof}
המאורע $E_i$ הוא פונקציה של $V$ והוא גורר $T_i=1$. נגדיר $Z_i=\mathbf 1_{E_i}$. אז $\Prb[T_i=0\mid Z_i=0]=1/[4(1-q_i)]$, וכן $H(T_i\mid Z_i=1)=0$ כל אימת שלמאורע זה מסה חיובית. מכאן $I(T_i;Z_i)=\phi(q_i)$. אי־שוויון עיבוד הנתונים נותן $\phi(q_i)\le I(B_i;V)$, ומ־\eqref{eq:budget} מתקבל האינדקס המבוקש. כדי שתתקיים ההבטחה, גם החלופה שנבנתה חייבת לקיים $q_i\ge\varepsilon$, ומכאן התנאי ההכרחי.
\end{proof}
לבלוקים בלתי תלויים ולכל $\varepsilon>0$ קבוע בתחום האמור, החסם ההכרחי הוא $m\ge (N/2)\phi(\varepsilon)=\Omega(N)$; אי־תלות הבדיקות אינה נדרשת. החסם הלא ליניארי מבטיח פער עקב אי־הבחנה שאינו נעלם עבור $m=o(N)$ כאשר $\limsup\TC(B)/N<\tfrac12 h_2(1/4)\approx0.4056$. לכן המספר $0.1556$ הוא סף מספיק שמקורו באומדן הפשוט יותר, ולא נקודת מעבר חדה ואוניברסלית. גם הסף האנטרופי המשופר מספק תנאי מספיק בלבד, ולא סיווג שלם המבוסס על המתאם הכולל בלבד. טענה \ref{prop:alias} מספקת את הקריטריון המדויק לאימות. בבלוקים בעלי מתאם חזק עשוי להיות מספר קטן של דפוסי יעד: למשל, אם כל הבלוקים זהים לאותו בלוק אחיד, די בבדיקת \EN{OR} אחת כדי לאמת את כולם.

\subsection{משמעות התוצאה לביקורת של מודלי שפה גדולים}
המשפט חל על התפלגות משותפת נתונה של תצורות התנהגות בינאריות, ולא ישירות על משקלי המודל, על הנחיות שרירותיות או על מדיניות הערכה אקראית. כדי להחילו יש להגדיר את אוכלוסיית התצורות, פרופיל דטרמיניסטי קבוע (לרבות הערכים האתחוליים שהוקפאו לצורך האקראיות), את הבלוקים בני שתי הסיביות והתפלגויותיהם השוליות, וכן את משפחת החלופות הקבילות. אם $N_{\rm eff}$ מציין את מספר הסיביות ההתנהגותיות שנבחרו, הוא שווה ל־$N$ במודל שהוגדר, ואינו אומד חדש של הממד. חסם הכרחי על מספר הבדיקות נובע רק אם ההבטחה הנדרשת חלה על החלופות שבמשפט, או על משפחה שהוכח כי היא כוללת את הנציגים הרלוונטיים. קיומה של מידה ששוקללה מחדש אינו גורר קיום של מודל יחיד, כוונון עדין מעשי או מתאם \EN{LoRA} המייצרים מידה זו. גם פרדיקט ה־\EN{OR} המופשט כשלעצמו אינו קריטריון בטיחות.

\section{עמימות סופית אינה אי־הבחנה אחידה}
יהי $X$ מרחב הסדרות הבינאריות הקבועות החל ממקום מסוים, עם טופולוגיית תת־המרחב המושרית מן המכפלה $\bits^{\mathbb N}$ שבה ערכי הקואורדינטות בדידים. נגדיר $X_j=\{x:\lim_i x_i=j\}$ וכן $T_i(x)=x_i$. הן $X_0$ והן $X_1$ צפופות ב־$X$: אפשר להרחיב כל קבוצת קואורדינטות סופית שערכיה נקבעו מראש לסדרה עם כל אחד משני הזנבות הקבועים. לכן יעד הזנב $\theta(x)=\lim_i x_i$ אינו רציף בשום נקודה כהעתקה למרחב הבדיד בן שתי נקודות.

מכך לא נובע שקיימות שתי מידות הסתברות קבועות שמסתירות את ההבדל מפני כל הבדיקות הסופיות. נבחן מידות הסתברות סיגמא־אדיטיביות $\mu$ על $X_1$ ו־$\nu$ על $X_0$, וכן בודק הבוחר מתוך $\mathcal T=\{T_i,1-T_i:i\in\mathbb N\}\cup\{0\}$. נגדיר $g(T,\mu,\nu)=\mathbb E_\mu T-\mathbb E_\nu T$, ו־$\lambda$ תנוע על פני כל התערובות הסיגמא־אדיטיביות של בדיקות.
\begin{theorem}[פער מינימקס מרבי]\label{thm:mm}
\[
\sup_\lambda\inf_{\mu,\nu}\mathbb E_{T\sim\lambda}g(T,\mu,\nu)=0,
\qquad
\inf_{\mu,\nu}\sup_{T\in\mathcal T}g(T,\mu,\nu)=1.
\]
\end{theorem}
\begin{proof}
לכל $\lambda$ ולכל $\eta>0$, קבוצה סופית של בדיקות נושאת מסה של לפחות $1-\eta$. בוחרים סדרה אחת מכל מחלקת זנב כך ששתיהן מסכימות בכל הקואורדינטות הנבדקות. מידות דיראק שלהן נותנות תשלום שאינו עולה על $\eta$. בדיקת האפס הקבועה מבטיחה תשלום אפס; בכך מוכח השוויון הראשון. לכל זוג מידות קביל, ממשפט ההתכנסות הנשלטת מתקבלים $\mathbb E_\mu T_i\to1$ וכן $\mathbb E_\nu T_i\to0$. הסופרמום הוא אחת, וכך גם האינפימום שלו על פני הזוגות.
\end{proof}
הפער הוא בין שני משחקים מוגדרים, ולא בין שני אומדנים של אותו משחק סופי. מחלקות הזנב אינן סגורות במרחב המכפלה הקומפקטי הסובב אותן. בפרט, מחלקות המידות שתומכיהן מוכלים בהן אינן קומפקטיות שם בטופולוגיה החלשה: סדרה של מידות דיראק בנקודות שתחילת זנבן הקבוע נדחית יכולה להתכנס חלש למידת דיראק במחלקה הנגדית. על כן לא מתקיימות כאן הנחות הקומפקטיות העשויות להצדיק החלפת סדר המינימקס \cite{sion}. ההוכחה נשענת על סיגמא־אדיטיביות; אין להרחיבה לאסטרטגיות אדיטיביות סופית שרירותיות ללא נימוק נוסף. הדוגמה מבודדת את כישלון העיקרון המאפשר להסיק עֵד גלובלי מקיומם של עדים לכל משפחת בדיקות סופית; היא אינה מזהה את כלל החסמים בתורת הסיבוכיות עם אובדן מידע בתצפית.

\section{שינוי אופן הגישה: מעבדת אורקלים מולטיליניארית}
\subsection{המודל האסימטרי והדרגה המזערית}
לפי מוסכמת ההרחבה מעל שדות סופיים של אהרונסון וויגדרסון (\EN{AW}) \cite{aw}, אורקל $A$ מורכב מטבלאות בוליאניות $A_r:\bits^r\to\bits$. הרחבה $\Atil_{r,\F}$ תואמת את הטבלה על הקובייה הבוליאנית ודרגתה בכל משתנה היא לכל היותר $c$. כאשר הדרגה היא לכל היותר אחת, ההרחבה יחידה מעל כל שדה סופי. שאילתה מפורשת כוללת את השדה ואת הקואורדינטות; משתמשים בייצוג המקובל של איברי שדה, המאפשר חישובים בזמן פולינומי. הסימון $[\mathrm{poly}]$ מציין את מוסכמת \EN{AW} להגבלת אורך השאילתה וגודל השדה. נגדיר
\[
u(A,\Atil)=\mathbf1[\NP^A\subseteq\PP^{\Atil}],\qquad
v(A,\Atil)=\mathbf1[\NP^{\Atil}\subseteq\PP^A].
\]
מכיוון שהאורקל המורחב יכול לדמות שאילתות בוליאניות, מתקיים $v\le u$. הפרופילים היחידים האפשריים מבחינה לוגית הם $00,10,11$. משפטים 5.1 ו־5.11(i) בקדם־הפרסום המקוון של \EN{AW} \cite{aw} מספקים עדים מולטיליניאריים ל־$11$ ול־$00$, בהתאמה; העד האחרון מתקבל מהפרדה חזקה יותר מ־$\BB^{\Atil}$. כאן אנו נותנים עד ישיר לפרופיל הביניים. איננו טוענים שתוצאה זו מחזקת את תוצאות האי־אפשרות של \EN{AW}; גם שאלת הקדימות של הבנייה המסוימת לפרופיל הביניים לא נבדקה באופן ממצה.

\begin{proposition}[דרגה אפס]\label{prop:degree0}
אם מרחיבים את המחלקה המותרת כך שתכלול גם דרגה אפס, אז $u=v$ לכל זוג קביל.
\end{proposition}
\begin{proof}
כל טבלה $A_r$ חייבת להיות קבועה על הקובייה הבוליאנית שלה. מצב זה מתאר אורקל אונרי (\EN{tally oracle}), שאינו בהכרח חישובי או טריוויאלי. שאילתה אלגברית בעלת $r$ ארגומנטים מחזירה את הערך הקבוע, שאפשר לקבל גם באמצעות השאילתה הבוליאנית $0^r$. בכיוון ההפוך, שאילתה בוליאנית היא מקרה פרטי של שאילתה אלגברית. שתי ההדמיות מוסיפות תקורה פולינומית למחלקות הזמן הנדונות, ולכן שתי מחלקות \EN{P} שוות זו לזו, וכך גם שתי מחלקות \EN{NP}.
\end{proof}
השימוש ב־$c=0$ כאן נוגע לתת־המחלקה שדרגתה אפס. כאשר \EN{AW} דורשים חסם שהוא מספר שלם חיובי, זהו מקרה מנוון המצורף במפורש. אין בכך זיהוי של כלל הגישה הבוליאנית הרגילה לאורקל עם דרגה אפס.

\subsection{פעולת הסכימה הבסיסית}
ההרחבה המולטיליניארית היחידה מתקבלת באמצעות אינטרפולציית לגראנז' על הקובייה הבוליאנית. עבור $x\in\bits^r$ ו־$s$ ארגומנטים מסיימים השווים ל־$2\in\F_3$,
\begin{equation}\label{eq:aggregate}
2^s\,\Atil_{r+s,\F_3}(x,\underbrace{2,\ldots,2}_{s})
=\sum_{y\in\bits^s}A_{r+s}(x,y)\pmod3.
\end{equation}
אכן, כאשר ערכו של ארגומנט מסיים הוא $2$, שני הגורמים $z$ ו־$1-z$ שווים ל־$2$ ב־$\F_3$, בעוד שהרישא הבוליאנית בוחרת את השורה המתאימה; בשדה מתקיים $2^{2s}=1$. זוהי פעולת הסכימה המולטיליניארית הידועה \cite{jkrs,aw}. בשורה שמובטח כי היא מכילה לכל היותר ביט אחד שערכו שונה מאפס, די בשאילתה יחידה כדי להכריע אם השורה אינה ריקה; השאילתה אינה מחזירה את הכתובת של אותו ביט. אלגוריתם דטרמיניסטי המבוסס על שאילתות בוליאניות ומכריע זאת לכל שורה המקיימת את ההבטחה זקוק ל־$2^s$ שאילתות במקרה הגרוע ביותר: כל עוד לא נבדקה כתובת מסוימת, תמליל שכולו אפסים אינו שולל את האפשרות שבכתובת זו הערך הוא אחת. הסכימה אינה מוסיפה מידע בלתי תלוי לטבלה המלאה, אך משנה את עלות הגישה לפונקציה של אותה טבלה.

\subsection{עד מפורש לפרופיל הביניים}
\begin{theorem}[הפרדה אסימטרית מולטיליניארית]\label{thm:w10}
קיים אורקל בוליאני $A$ שהרחבתו המולטיליניארית היחידה מעל כל השדות הסופיים מקיימת
\begin{equation}\label{eq:w10}
\PP^A\subsetneq\UP^A=\NP^A\subseteq\PP^{\Atil[\mathrm{poly}]}\subseteq\PP^{\Atil}.
\end{equation}
בפרט, הפרופיל $10$ מתממש כאשר הדרגה בכל משתנה היא אחת, וחסם הדרגה המזערי המאפשר אותו (כאשר כוללים גם אפס) הוא אחת.
\end{theorem}
\begin{proof}
נשתמש במכונות אורקל סטנדרטיות עם חסם זמן וסרטי שאילתה ריקים בתחילת הריצה. כך ריצה בת $t$ צעדים מגישה לכל היותר $t$ שאילתות, שאורך כל אחת מהן אינו עולה על $t$. נגדיר את שפת הקבלה היחסית בזמן חסום $K^A$ באמצעות הקידוד בעל התיחום העצמי
\[
e(i,w,t)=0^i1\,0^{|w|}1\,w\,1^t,
\]
כאשר $i$ הוא אינדקס של מכונה אי־דטרמיניסטית בקידוד אפקטיבי סטנדרטי (למשל, מספרים טבעיים המקודדים תיאורים בינאריים של מכונות). מילה כזאת שייכת לשפה אם מכונה $i$ מקבלת את $w$ תוך לכל היותר $t$ צעדים; קלט שאינו בקידוד התקין נדחה. אורך המילה הוא $i+2|w|+t+2>t$, ולכן $K^A(x)$ עבור $|x|=r$ תלוי רק בתשובות האורקל באורכים הקטנים מ־$r$. השפה שייכת ל־$\NP^A$ ושלמה למחלקה זו ביחס לרדוקציות דטרמיניסטיות מסוג רבים־לאחד בזמן פולינומי: עבור מכונה קבועה עם חסם זמן, שולחים את $w$ אל קידוד המורכב מקוד המכונה, מ־$w$ ומחסם הזמן הפולינומי שלה. הדמיה אוניברסלית מוסיפה תקורה פולינומית באורך הקלט המקודד.

נבנה את $A$ על ידי קביעת רישאות סופיות הולכות וגדלות לפי אורך. באורכים זוגיים חיוביים $2r$, כאשר $r\ge1$, נדרוש לכל $x\in\bits^r$ כי
\begin{equation}\label{eq:gi1}
\sum_{y\in\bits^r}A(xy)=\mathbf1_{K^A}(x),
\end{equation}
כאשר הסכום הוא על המספרים השלמים. בכל אורך אי־זוגי נדרוש לכל היותר ביט אחד שערכו אחת. ערכה של המילה הריקה נקבע לאפס. שני האינווריאנטים נשמרים בכל שלב, וכך גם ההסכמה עם כל האורכים שכבר נקבעו.

נמנה כל זוג $(M_i,p_i)$ של מכונת אורקל דטרמיניסטית וחסם זמן פולינומי ממשפחה בת־מנייה וקופינלית של חסמים, למשל $p_i(n)=(n+2)^k$. נניח שכל האורכים עד $B$ כבר נקבעו. נבחר $n>B$ אי־זוגי, עם $n\ge3$, כך ש־$2^{\lfloor n/2\rfloor}>p_i(n)$. נמלא לפי הסדר את כל האורכים שביניהם: באורכים אי־זוגיים נציב אפסים, ובכל אורך זוגי $2r$ נאפס שורה כאשר $K^A(x)=0$, ובמקרה האחר נציב אחת יחידה בכתובת הראשונה בסדר לקסיקוגרפי. הגדרת הקבלה בזמן חסום נעזרת רק באורכים קצרים יותר שכבר נקבעו.

נריץ את $M_i$ על $1^n$ במשך הזמן המוקצב, תוך שימוש ברישא הקבועה ומתן תשובת אפס לכל שאילתה שטרם נקבע לה ערך. נסמן את תשובת המכונה ב־$b$ (אם הזמן אזל, נחשיב זאת כדחייה), ונסמן ב־$Q$ את אוסף השאילתות שטרם נקבע להן ערך. נקבע $B'=\max(n,p_i(n))+1$. באורך $n$, נציב את הביט $1-b$ בכתובת הראשונה שלא נשאלה, ואת כל שאר הערכים נאפס. כתובת כזאת קיימת, שכן $|Q|\le p_i(n)<2^{\lfloor n/2\rfloor}<2^n$. באורכים האי־זוגיים האחרים עד $B'$ נציב אפסים. בכל אורך זוגי $2r$ בתחום שנותר נמלא את השורות בהתאם ל־\eqref{eq:gi1}, תוך שימוש בכתובת שלא נשאלה כאשר יש להציב אחת: הרי בכל שורה יש $2^r>|Q|$ כתובות, מאחר ש־$2r>n$. כל הקביעות נעשות לפי סדר אורך עולה, ולכן ערכי $K^A(x)$ מוגדרים בעת הקביעה. לכל שאילתה בריצה שתועדה נשמרת בדיוק התשובה הקודמת; הדטרמיניזם מבטיח שאותה ריצה תתקבל שוב, והשלבים הבאים לעולם אינם משנים אורכים שכבר נקבעו. בכך מסתיים השלב.

איחוד הרישאות העולות מגדיר את $A$ בכל אורך. אפשר לבצע את כל הבחירות באמצעות הכתובות הסופיות הקטנות ביותר והאורכים הקבילים הקטנים ביותר, ולכן הבנייה אינה דורשת את אקסיומת הבחירה. נגדיר $L=\{1^n:n\text{ odd},\ \exists y\in\bits^n\ A(y)=1\}$. השלב של כל מכונה בעלת חסם זמן מאלץ אותה להשיב בניגוד לשייכות ל־$L$. לכן $L\notin\PP^A$. אינווריאנט היחידות באורכים האי־זוגיים נותן $L\in\UP^A$ באמצעות ניחוש $y$ ושאילתת $A(y)$; יש לכל היותר ענף מקבל אחד. משוואה \eqref{eq:gi1} נותנת באותו אופן $K^A\in\UP^A$. סגירות תחת הרדוקציות הדטרמיניסטיות הקבועות נותנת $\NP^A=\UP^A$. לפי \eqref{eq:aggregate}, אפשר להכריע את $K^A$ בשאילתה יחידה בנקודה $(x,2,\ldots,2)$ מעל $\F_3$, בעלת $2|x|$ ארגומנטים, ולאחריה בדיקה אם התוצאה שונה מאפס (בקלט ריק, שאינו בקידוד התקין, מטפלים ישירות). יחד עם השלמות מתקבל $\NP^A\subseteq\PP^{\Atil[\mathrm{poly}]}$. מכיוון ש־$\PP^A\subseteq\UP^A$ ו־$L$ מפרידה ביניהן, מתקבלת \eqref{eq:w10}. הפרופיל הוא $10$: מתקיים $u=1$ מכוח ההכלה, ו־$v=0$ משום ש־$\NP^{\Atil}\supseteq\NP^A\ni L\notin\PP^A$. טענה \ref{prop:degree0} שוללת דרגה אפס.
\end{proof}
אפשר להשתמש באותו עד לפרופיל הביניים גם עם ההרחבה המולטיליניארית מעל השלמים: צמצום תוצאת החישוב השלמה מודולו שלוש נותן את \eqref{eq:aggregate}, ומספר הביטים של פלט ההרחבה מעל השלמים פולינומי: הערך המוחלט של סכום האינטרפולציה, הכולל לכל היותר $2^r$ מכפלות עבור $r$ ארגומנטים, אינו עולה על $2^r\prod_j(|z_j|+1)$. התוצאה אינה מראה שמכונה בזמן פולינומי ללא אורקל יכולה לחשב את הסכימה מתיאור תמציתי, או ששתי המחלקות המורחבות \EN{P} ו־\EN{NP} שוות.

\Needspace{9\baselineskip}
\subsection{כל תצפית סופית מותירה שלושה פרופילים אפשריים}
\begin{theorem}[עמימות מתמשכת של משתנה המטרה]\label{thm:top}
יהי $\mathcal X$ מרחב האורקלים הבוליאניים עם הרחבותיהם המולטיליניאריות היחידות, המצויד בטופולוגיה של תצפיות רגילות ואלגבריות סופיות. כל גליל תצפיות סופי שאינו ריק מכיל עולמות מכל שלושת הפרופילים $00,10,11$. לכן העתקת הפרופיל לקבוצה הדיסקרטית בת שלוש הנקודות אינה רציפה באף נקודה.
\end{theorem}
\begin{proof}
אם שני אורקלים בוליאניים מסכימים בכל האורכים שמעל $M$, אז הרחבותיהם המולטיליניאריות היחידות מסכימות באורכים אלה מעל כל שדה. במספר הסופי של אורכי הקלט הנותרים, טבלת אמת סופית המקודדת מראש ואינטרפולציה מדמות כל שאילתת הרחבה בתקורה פולינומית באורך השאילתה המפורשת: מספר האיברים בסכום חסום על ידי $2^M$, שהוא קבוע, ופעולות החשבון בשדה יעילות בייצוג המפורש. אפשר לדמות גם שאילתות רגילות. מכאן ששינויים סופיים משמרים כל אחת ממחלקות \EN{P} ו־\EN{NP} ואת הפרופיל האסימטרי, לרבות תחת מוסכמת ההרחבה החסומה. יהי נתון תמליל סופי של תצפיות שניתן לממש, ונבחר אורקל אחד שמממש אותו. אפשר להחליף אצל כל אחד משלושת העדים הגלובליים (משפט \ref{thm:w10} ושתי תוצאות הקצה של \EN{AW}) את הטבלאות עד לאורך המרבי שהופיע בשאילתות בטבלאות של האורקל המממש. יחידות ההרחבה משמרת כל ערך שנצפה, ואי־התלות בשינויים סופיים משמרת את הפרופיל של כל עד. לפיכך כל סביבה פוגשת את כל שלוש מחלקות הפרופיל; רציפות לתוך קבוצה דיסקרטית בלתי אפשרית בכל נקודה.
\end{proof}
המשפט עוסק בטופולוגיה של עולמות אורקל, ולא בטופולוגיה של שיטות הוכחה. הוא גם ממחיש מדוע פעולת הסכימה השימושית בשאילתה יחידה ב־\eqref{eq:aggregate} אינה חושפת פרופיל סיבוכיות גלובלי באמצעות אוסף סופי קבוע של תצפיות. הפרדות יחסיות לאורקל מתיישבות עם כל אחת משתי התשובות האפשריות לבעיית \EN{P versus NP} ללא אורקל.

\section{אימות ממוחשב: מה בדיוק נבדק}
החבילה שסופקה כוללת שישה קובצי \EN{Lean 4.36.0} עצמאיים ללא ייבוא חיצוני מפורש.

הקובץ \EN{\texttt{AXIOM\_AUDIT.log}} שסופק מתעד 40 דוחות ברמת ההצהרה (ארבעה ללא אקסיומות, וביתר רק \EN{\texttt{propext}} ו־\EN{\texttt{Quot.sound}}).

הקובץ \EN{\texttt{RELEASE\_GATE.log}} מתעד קומפילציה נפרדת מוצלחת ו־32 בדיקות של הצהרות שורש; הרישום המורחב של 40 הדוחות הוא בדיקה נפרדת. התחיליות של \EN{SHA-256} להלן מזהות את גרסאות קובצי המקור שסופקו, ואינן מעידות על הידור חדש לצורך קדם־פרסום זה.
\begin{center}\small
\begin{tabular}{@{}llrl@{}}
\toprule
\textbf{מודול} & \textbf{תחילית \EN{SHA-256}} & \textbf{דוחות} & \RTLcell{0.43\linewidth}{\textbf{היקף}}\\[4pt]
\midrule
\EN{\texttt{F3Multilinear}} & \EN{\texttt{eb37a377aa7d}} & 4 & \RTLcell{0.43\linewidth}{התאמה בוליאנית, יחידות ההרחבה המולטילינארית, פענוח שורה}\\[4pt]
\EN{\texttt{HiddenRow}} & \EN{\texttt{8d80d618f611}} & 5 & \RTLcell{0.43\linewidth}{כתובות פנויות, עדכונים נקודתיים, יחידות}\\[4pt]
\EN{\texttt{AcyclicKernel}} & \EN{\texttt{f0d873487b41}} & 5 & \RTLcell{0.43\linewidth}{קידוד, פענוח תחבירי, תקציב צעדים חסום ממש, מקומיות}\\[4pt]
\EN{\texttt{OracleDiagonal}} & \EN{\texttt{7bff2ef97995}} & 5 & \RTLcell{0.43\linewidth}{גבול אינסופי, התאמה בין שלבים, אינווריאנטי \EN{GI}, ואלכסון ביחס למודל המכונה המופשט}\\[4pt]
\EN{\texttt{AsymmetricCollapse}} & \EN{\texttt{3e9c83638a54}} & 8 & \RTLcell{0.43\linewidth}{קומפוזיציה מותנית החל מ־\EN{\texttt{VerifiedLevel01}}}\\[4pt]
\EN{\texttt{ProfileTopology}} & \EN{\texttt{f9cf22a3a9ca}} & 13 & \RTLcell{0.43\linewidth}{סדר הפרופילים, לוגיקה של דרגה אפס, עמימות זנב מותנית, עדי זנב בדידים}\\[4pt]
\bottomrule
\end{tabular}
\end{center}
הלוגים מדווחים על אפס שגיאות ואזהרות, ועל אפס מופעים של \EN{\texttt{sorryAx}}, של \EN{\texttt{ofReduceBool}} ושל \EN{\texttt{Classical.choice}}. זו ראיית הידור שסופקה, לא קומפילציה עצמאית שבוצעה לגרסה זו. ארבעת המודולים הראשונים מצהירים מחדש על שמות משותפים, כך שאי אפשר לייבא או לשרשר אותם לספרייה אחת ללא שינויים. בפרט, \EN{\texttt{asymmetric\_collapse\_root}} מקבל ארגומנט מסוג \EN{\texttt{VerifiedLevel01}}; אין בנאי שנבדק ומקשר ארגומנט זה לקבצים האחרים. גם התאמת מודל מכונת האורקל למודל הסטנדרטי, והגשר המלא ברמת המחלקות בין \EN{UP/NP} לבין זמן אלגברי, לא הושלמו פורמלית. מודול טופולוגיית הפרופילים מניח את קיום העדים ואת האינווריאנטיות לשינויים סופיים. לפיכך משפטים \ref{thm:w10} ו־\ref{thm:top} לעיל הם טענות אנליטיות הנתמכות ברכיבים מאומתים, לא משפטים בלתי־מותנים שאומתו מקצה לקצה ב־\EN{Lean}. הוכחות האנטרופיה וההסתברות בסעיפים \EN{2--3} לא אומתו ב־\EN{Lean}.

לשם שחזור באמצעות החבילה המקורית, יש לבדוק את המניפסט המלא באמצעות \EN{\texttt{shasum -a 256 -c SHA256SUMS.txt}} ולהריץ \EN{\texttt{bash verify\_release.sh}} בתיקיית ההפצה באמצעות \EN{Lean 4.36.0}. שער הבדיקה הנוכחי אינו אוכף גרסה מדויקת של כלי הבנייה, אינו פוסל כל אזהרה, אינו אוכף רשימת אקסיומות מותרות המכסה את כל האפשרויות, ואינו בודק את כל 40 השורשים או את הרכבתם. אלו דרישות מפורשות להנדסת ההפצה, שסיכום מספרי של הבדיקות לבדו אינו ממלא. קובץ ה־\EN{README} שסופק מתעד השפעות מקומיות של השימוש ב־\EN{\texttt{simp}} וב־\EN{\texttt{omega}} על התלות באקסיומות; אלה תצפיות על ההוכחות המסוימות, לא טענות כלליות על הטקטיקות.

\section{פרוטוקול אמפירי מוצע (לא בוצע)}
בדיקת התחולה על מודלי שפה גדולים חייבת להפריד בין אי־הבחנה בין פרופילים במובן המתמטי לבין יכולת המימוש במשפחת מודלים. להלן פרוטוקול מוצע; אין כאן טענה לרישום מוקדם חיצוני הנושא חותמת זמן, לאיסוף מדגם או לביצוע מדידה.
\begin{enumerate}
\item קבעו מראש משפחת מודלים, משימות התנהגותיות לא מזיקות, בלוקי יעד, זרעי אקראיות, מערך ביקורת וקבוצת הערכה נפרדת, לפני צפייה בתוצאות. הגדירו קונפיגורציה כרשומה מלאה של התשובות בפרוטוקול המוגדר. בדקו אם ההתפלגויות השוליות של הבלוקים קרובות לאחידות; אם לא, התאימו במפורש את המודל האנליטי או דחו את ההחלה הישירה של המשפט.
\item אִמדו את $H(B)$ ואת האנטרופיות של הבלוקים לפי מודל דגימה משותף ומתועד. השתמשו בחסמי אי־ודאות או באומדים מבניים שנבדקו בהתפלגויות סינתטיות בעלות מתאם כולל ידוע; בדיקות אי־תלות שוליות לבדן אינן קובעות שהמתאם הכולל נמוך. דווחו את הטווח שהתקבל עבור $m/N$ ועבור $\TC(B)/N$, במקום להתייחס ל־$N_{\rm eff}$ כאל קבוע שנקבע אמפירית.
\item במודל סינתטי סופי חשבו את האופטימום המדויק המשמר את הפרופיל, כבטענה \ref{prop:alias}, כביקורת חיובית. בניסוי נפרד ובעל תקציב מוגדר להתאמת מודלים, בדקו אם מודל חלופי קביל משחזר את חוק פרופיל הביקורת תוך שינוי יעד לא מזיק שהוגדר מראש ונשמר לבדיקה. הפרידו את נתוני הביקורת מנתוני ההערכה השמורים והשוו לבקרות ללא שינוי ולבקרות שהותאמו באקראי.
\item הגדירו מראש סף מרחק מותר בין ההתפלגויות המשותפות וסף להשפעה על היעד, והציגו רווחי סמך. שוויון בין ממוצעי הבדיקות הבודדות אינו ראיה לשוויון החוק המשותף של הפרופיל; אמידת החוק המשותף בממד גבוה עלולה כשלעצמה להיות בלתי ישימה. כישלון במימוש הוא ראיה לגבי המשפחה ושיטת האופטימיזציה שנבחרו, לא הפרכה של המשפט הסופי.
\end{enumerate}
הבחירה המעשית בין הוספת בדיקות, פעולת גישה חדשה או הגבלה של משפחת המודלים מחייבת מודל נפרד של עלות ונכונות. גישה לקופסה לבנה לא תמיד מספיקה, וביקורת סופית של קופסה שחורה אינה תמיד חסרת תועלת.

\section{סיכום}
שימור מדויק של פרופיל, מעבר מתאימות סופית לתאימות גלובלית וגישה חישובית הם עניינים נבדלים. בביקורת סופית, תנאי הפירוק המדויק והאומדנים התלויים במתאם מבהירים מה ניתן לאמת עבור תומך מוגדר ומודל נתון של מידות חלופיות. דוגמת הזנב שוללת החלפה בלתי מוגבלת של סדר הכמתים על העדים ועל הבדיקות. בניית האורקל המולטילינארי מבססת באופן מפורש משטר ביניים אסימטרי בחסם הדרגה המזערי, וממחישה את כוחה של פעולת גישה מאגדת גם כשהיא נקבעת מאותם נתונים בסיסיים. החוט המקשר הוא שאלת מחקר משותפת, לא רדוקציה של בטיחות מודלי שפה או של בעיית \EN{P} מול \EN{NP} לאי־שוויון האנטרופיה הסופי. הצעד הפורמלי הבא הוא הוכחת שורש משולבת הכוללת גישור למכונה הסטנדרטית; הצעד האמפירי הבא הוא לקבוע אם ההנחות ההתנהגותיות ויכולת המימוש אכן מתקיימות במשפחת מודלים ממשית.

\begin{english}
\renewcommand{\refname}{References}
\begin{thebibliography}{9}\small
\bibitem{cook} S. Cook. \emph{The P versus NP Problem}. Clay Mathematics Institute, 2000. \url{https://www.claymath.org/wp-content/uploads/2022/06/pvsnp.pdf}.
\bibitem{aw} S. Aaronson and A. Wigderson. Algebrization: A New Barrier in Complexity Theory. \emph{ACM Transactions on Computation Theory}, 1(1), 2009. Theorem numbering follows the online preprint: \url{https://www.scottaaronson.com/papers/alg.pdf}.
\bibitem{bgs} T. Baker, J. Gill, and R. Solovay. Relativizations of the P =? NP question. \emph{SIAM Journal on Computing}, 4(4):431--442, 1975.
\bibitem{jkrs} A. Juma, V. Kabanets, C. Rackoff, and A. Shpilka. The Black-Box Query Complexity of Polynomial Summation. \emph{Computational Complexity}, 18(1):59--79, 2009. \url{https://eccc.weizmann.ac.il/report/2007/125/}.
\bibitem{rr} A. A. Razborov and S. Rudich. Natural Proofs. \emph{Journal of Computer and System Sciences}, 55(1):24--35, 1997.
\bibitem{cover} T. M. Cover and J. A. Thomas. \emph{Elements of Information Theory}. 2nd ed., Wiley, 2006.
\bibitem{watanabe} S. Watanabe. Information theoretical analysis of multivariate correlation. \emph{IBM Journal of Research and Development}, 4(1):66--82, 1960.
\bibitem{sion} M. Sion. On general minimax theorems. \emph{Pacific Journal of Mathematics}, 8(1):171--176, 1958.
\bibitem{taylor} J. Taylor et al. Auditing Games for Sandbagging. arXiv:2512.07810, 2025. \url{https://arxiv.org/abs/2512.07810}.
\end{thebibliography}
\end{english}
\end{document}
